% Stage C1: the plain sieve in the skeleton of the C code, with the theorem that it returns the list
% of the brute-force reference. Proposed names are set with \lname and become \lean once the
% declarations exist.
\chapter{The plain sieve}\label{chap:sieve}

Stage~C1 replaces the brute-force reference in the program by a sieve in the skeleton of the C
function \verb|find_points_work_1|, with the simplest component in each place: the denominators
$1, \dots, H$ in order, the numerators $-H, \dots, H$, one Boolean per numerator, prime moduli
with tables of admissible residues, the test for common factors, the exact check, the delivery of
$\pm y$, and the points at infinity computed directly from $d$ and $c_d$. Every row of the sieve is
\emph{sound}: it lets every point through. So for every well-formed choice of the sieving
parameters the sieve returns exactly the list of the brute-force reference, the same points in the
same order, and membership and no duplicates follow from Theorem~\ref{thm:bruteForce}.

The later stages contribute in four ways. A \emph{data refinement} computes the same predicate in
another way, with a lemma that it does: the bit arrays of words of stage~C2, the faithful tables
and the packing of stage~F; the equation with the reference survives it. A \emph{transformation},
the reversal of stage~D, is a bijection of the point sets that preserves the box. A \emph{sound
filter} comes with the lemma that every point passes it: the 2-adic information and the tests on
the denominators of stage~E, the composite moduli and the third stage of stage~F. An
\emph{enumeration refinement} visits the denominators that carry points in another order: the loop
shapes of stage~E. From the reversal on, and with the divisor order of stage~E, the statement is
that the output is a permutation of the reference, which still gives membership and no duplicates;
stage~G proves the order of the program.

\section{What a default run of the C code does}\label{sec:skeleton}

\verb|find_points| (\csource{find_points.c}{5226}{5242}) calls \verb|find_points_work|, which
refuses a degree outside $0, \dots, \mathrm{maxDegree}$ (\csource{find_points.c}{4557}{4558}) and
calls
\verb|find_points_work_1|. In a run without options, that function:
\begin{enumerate}
  \item drops the leading zero coefficients, refusing the form as not squarefree when that lowers
    $\lceil n/2 \rceil$, refuses degree $0$ and a height below $1$, and caps the parameters
    (\csource{find_points.c}{4624}{4658}); it reverses the polynomial in some cases
    (\csource{find_points.c}{4732}{4803}; stage~D);
  \item computes the search region by the Sturm step (stage~H) and decides whether there are
    points at infinity (\csource{find_points.c}{4894}{4899});
  \item sets up the tests on the denominators (squares or divisors times squares, Jacobi symbols)
    and the 2-adic information, with an exit when no class of denominators modulo $64$ admits a
    numerator (stage~E);
  \item chooses the moduli and the forbidden divisors, examining each prime: the table
    \verb|is_f_square| of the residues $x$ modulo $p$ for which $f(x)$ is a square modulo $p$
    (\csource{find_points.c}{1703}{1747}), with an exit when there are no points modulo $p$
    (stage~D);
  \item delivers the points at infinity (\csource{find_points.c}{5068}{5105});
  \item runs over the denominators in blocks (\csource{find_points.c}{5152}{5159}); for each
    denominator $b$ that passes its tests, it sieves the numerators of the search region in pieces
    of bit arrays with the rows of the moduli of the first two stages (\verb|sift| and
    \verb|_ratpoints_sift0|), tests each survivor for common factors with $b$
    (\csource{sift.c}{318}{343}), then with the primes of the third stage one numerator at a time,
    and checks the rest exactly, delivering $(a, y, b)$ and, for $y \ne 0$, $(a, -y, b)$
    (\csource{find_points.c}{5262}{5326}).
\end{enumerate}
C1 models what decides which points are found, each in its simplest form. Most of the rest only
spares work; the refusals, the reversal and the exit when there are no points modulo $p$ (D) and
the options of the output (G) come in their stages.

\section{The parameters}\label{sec:params}

\begin{definition}[The primes of the sieve]\label{def:primeTable}
  \lean{Ratpoints.primeTable, Ratpoints.maxBitsInPrime, Ratpoints.numPrimes, Ratpoints.prime_of_mem_primeTable}
  \leanok
  The 171 odd primes below $1024$, in increasing order, the table of \verb|primes.h|
  (\csource{primes.h}{29}{40}); every entry is a prime. The moduli are chosen among those below
  $2^b$, for the size in bits $b = \mathrm{maxBitsInPrime}$ that the build of ratpoints fixes
  (\csource{rp-private.h}{321}{323}): $8$ in the default build, whose Makefile sets it, where
  the header takes $7$ when it is given none. There are $\mathrm{numPrimes}$ of them, $53$ for
  $b = 8$ (\csource{rp-private.h}{325}{361}, where the header lists the numbers for
  $b = 5, \dots, 10$).
\end{definition}

\begin{definition}[Sieving parameters]\label{def:Params}
  \uses{def:primeTable}
  \lean{Ratpoints.Params, Ratpoints.Params.WF}
  \leanok
  The free choices of the sieve: the moduli, a list of primes in the order of the stages; the cut
  points $s_1 \le s_2$, positions in that list, which end the first and the second stage (the
  third stage is the rest of the list); at most how many forbidden divisors are tested; and
  whether the Jacobi test runs. The parameters are \emph{well-formed} when every modulus is a
  prime of the table and $s_1 \le s_2 \le$ the number of moduli.
\end{definition}

\begin{definition}[The parameters of the executable]\label{def:sieveParams}
  \uses{def:Params, def:primeTable}
  \lean{Ratpoints.SieveOptions, Ratpoints.sieveParams, Ratpoints.defaultParams, Ratpoints.defaultNumPrimes, Ratpoints.defaultMaxForbidden, Ratpoints.wf_sieveParams, Ratpoints.wf_defaultParams}
  \leanok
  The parameters made from the values of the options \verb|-p|, \verb|-n|, \verb|-N|, \verb|-P|,
  \verb|-F| and \verb|-j|/\verb|-J|, with the defaults and caps of ratpoints
  (\csource{ratpoints.h}{333}{340}, \csource{ratpoints.h}{345}{352},
  \csource{find_points.c}{4643}{4649}, \csource{find_points.c}{4655}{4656},
  \csource{find_points.c}{4662}{4665}): $q = \min(k, \mathrm{numPrimes})$ primes are taken, for
  the value $k \ge 0$ of \verb|-p| or else $k = \mathrm{defaultNumPrimes} = 30$;
  $s_2 = \min(N, q)$ and $s_1 = \min(n, s_2)$ for the values of \verb|-N| and \verb|-n|,
  $s_2 = q$ and $s_1 = s_2$ when they are absent or negative; the moduli are the first
  $\min(q + P, \mathrm{numPrimes})$ primes of the table for the value $P \ge 0$ of \verb|-P|,
  the first $q$ when it is absent or negative; at most $\min(F, 171)$ forbidden divisors, for
  the value $F \ge 0$ of \verb|-F| or else $F = \mathrm{defaultMaxForbidden} = 64$; the Jacobi
  test unless \verb|-j| is the last of \verb|-j| and \verb|-J|. The default parameters are those
  of no options: the first 30 primes, all in the first stage. Every parameter set made this way
  is well-formed.
\end{definition}

\verb|Params| holds the choices that the theorems quantify over: which moduli the sieve uses, in
which order and in which stage. ratpoints makes them from the curve (\verb|sieving_info|, which
ranks the primes and may look past the first $q$ when the curve needs it), and moves the end of the
third stage during the run (\verb|adapt_primes|; with \verb|-A 2| also $s_2$). None of this
changes the points found or their order, since a modulus only removes candidates that are not
points (\verb|-q| suppresses the report of the moduli used, which does depend on the choice). So the
model takes the choice as an input, and the theorems hold for every well-formed choice. They
quantify over the parameters that \verb|Params.WF| admits, a definition of a few lines in the
model, which has no Mathlib, so the comparator challenge imports it; the runs of the executable
need no such reading, since \verb|wf_sieveParams|, certified with the main theorems, shows that
every parameter set the command line produces is well-formed, which also rules out a
well-formedness that nothing satisfies. C1 reads only the moduli and treats all stages alike; the
cut points, the forbidden divisors and the Jacobi switch are read from stages~E and~F on, and
further choices (the block length, the length of the pieces) join \verb|Params| with the
components that use them. Well-formedness constrains a field only as far as its component needs
it to be meaningful. The constants of the build that these definitions read, the size in bits of
the moduli and the defaults of \verb|-p| and \verb|-F|, are definitions of the model like
$\mathrm{maxDegree}$ (Definition~\ref{def:maxDegree}), and no proof depends on their values:
\verb|wf_sieveParams| needs only that $\mathrm{numPrimes}$, a number of primes of the table, is
at most its size.

\section{The table of a prime}\label{sec:primeTables}

Write $f = \sum_{i \le d} c_i x^i$ for the polynomial of degree $d$. For a prime $p$ and a
denominator $b$ not divisible by $p$, $F_c(a, b) \equiv b^D f(a b^{-1}) \pmod p$, and $b^D$ is the
square of a unit since $D$ is even; so if $F_c(a, b) = y^2$, then $f(a b^{-1})$ is a square
modulo $p$ (zero included). If $p > 1$ divides $b$ and $\gcd(a, b) = 1$, then $p$ does not divide
$a$. These two facts make sound the rows that the tables below define.

\begin{definition}[Squares and inverses modulo a prime]\label{def:squaresMod}
  Proposed names: \lname{Ratpoints.Sieve.squaresMod}, \lname{Ratpoints.Sieve.inversesMod},
  \lname{Ratpoints.Sieve.getElem!\_squaresMod\_mul\_self\_mod},
  \lname{Ratpoints.Sieve.getElem!\_inversesMod\_mul\_mod}.
  The table of the squares modulo $p$, made as the generator of the tables of ratpoints makes it:
  zero, and the squares of the odd residues (\csource{gen_find_points_h.c}{69}{88}); and the
  table of the inverses modulo $p$ of the residues $1, \dots, p - 1$, each the least positive one
  found by trial (\csource{gen_find_points_h.c}{33}{44}, \csource{gen_find_points_h.c}{90}{100}).
  The squares of all residues are marked, and $r \cdot r^{-1} \equiv 1 \pmod p$ for a prime $p$.
\end{definition}

\begin{definition}[The table of a prime]\label{def:examinePrime}
  \uses{def:squaresMod}
  Proposed names: \lname{Ratpoints.Sieve.evalMod}, \lname{Ratpoints.Sieve.Entry},
  \lname{Ratpoints.Sieve.examinePrime}, \lname{Ratpoints.Sieve.natCast\_evalMod}.
  For a prime $p$: the coefficients reduced modulo $p$, and for each residue $x$ modulo $p$
  whether $f(x)$ is a square modulo $p$, with $f(x)$ computed by Horner's rule modulo $p$
  (\csource{find_points.c}{1703}{1747}); kept, as the sieve entry of ratpoints keeps it
  (\csource{rp-private.h}{542}{555}), together with $p$ and the inverses modulo $p$
  (\csource{find_points.c}{1795}{1824}).
\end{definition}

\begin{definition}[The rows of a denominator]\label{def:admits}
  \uses{def:examinePrime}
  Proposed names: \lname{Ratpoints.Sieve.Entry.admits}, \lname{Ratpoints.Sieve.survives}.
  The row of the denominator $b$ in the table of $p$ admits the numerator $a$ when, for the
  residues $r = b \bmod p$ and $x = a \bmod p$, the table's entry at $x r^{-1} \bmod p$ is true if
  $r \ne 0$, and $x \ne 0$ if $r = 0$. This is what the bit of $a$ says in the row of $r$ in
  ratpoints, where the rows are bit patterns built from the table (\csource{init.c}{118}{246})
  and, for $r = 0$, a fixed pattern that does not depend on the curve
  (\csource{gen_find_points_h.c}{102}{134}). A numerator \emph{survives} the moduli for $b$ when
  the row of $b$ in the table of every modulus admits it; this test, a function of the numerator,
  is the interface between the sieve and the check, and the bit arrays of stage~C2 are specified
  against it.
\end{definition}

\begin{lemma}[A square value of the form gives a square value of $f$]\label{lem:isSquare_div}
  Proposed name: \lname{Ratpoints.isSquare\_sum\_mul\_div\_pow\_of\_sq\_eq}.
  In a field, if $b \ne 0$, $D$ is even, $d \le D$ and
  $y^2 = \sum_{i \le d} f_i\, a^i b^{D - i}$, then $\sum_{i \le d} f_i\, (a/b)^i$ is the square of
  $y / b^{D/2}$.
\end{lemma}

Mathlib has the identity behind this in the form of polynomials, \verb|Polynomial.eval_homogenize|;
the sum form, with its few lines of proof, suits the definition of $F_c$.

\begin{lemma}[The rows let every point through]\label{lem:admits_examinePrime}
  \uses{def:admits, def:binaryForm}
  Proposed name: \lname{Ratpoints.Sieve.admits\_examinePrime\_of\_sq\_eq}.
  Let $p$ be a prime, $b$ a natural number, $\gcd(a, b) = 1$ and $y^2 = F_c(a, b)$. Then the row
  of $b$ in the table of $p$ for the degree of $c$ admits $a$.
\end{lemma}

\begin{proof}
  \uses{lem:isSquare_div}
  The tables are what they claim to be: the squares, the inverses and Horner's rule modulo $p$
  compute their definitions. For $p \nmid b$, cast $y^2 = F_c(a, b)$ to $\ZZ/p\ZZ$ and apply
  Lemma~\ref{lem:isSquare_div}; the table of squares then has the entry at $a b^{-1}$. For
  $p \mid b$, a common divisor $p > 1$ of $a$ and $b$ would divide $\gcd(a, b) = 1$, so
  $p \nmid a$.
\end{proof}

Lemma~\ref{lem:admits_examinePrime} needs neither $b \ge 1$ nor an odd prime, so at its use the
sieve's denominators and the primes of the table need no side conditions.

\section{The sieve}\label{sec:sieve}

The model follows the rules of the translation throughout. A C \verb|long| is an \verb|Int|; the
obligation that its value lies in the range of the type comes with the faithful components that
compute in \verb|long|s (D, F). The remainder \verb|mod()| of the C code is Lean's Euclidean
\verb|%|. The tables that the build of ratpoints generates (the squares, the inverses, the row of
$0$) are Lean functions with lemmas that they compute what they claim, checked against the
generated header by \verb|#guard| examples, and as a whole with the faithful rows of stage~F. Each
loop of the sieve is a definition of type \verb|Id α| that its callers bind with \verb|←|, so that
its contract, an \verb|@[spec]| theorem proved with the program logic of Lean core and with loop
invariants that are named definitions, applies at the call; only \verb|findPoints| runs a
component with \verb|Id.run|. The tables of a prime are folds and maps, specified by ordinary
lemmas.

\begin{definition}[The sieve of one denominator]\label{def:sieveDen}
  \uses{def:admits}
  Proposed names: \lname{Ratpoints.Sieve.sieveDen}, \lname{Ratpoints.Sieve.sieveDen\_spec}.
  An array of $2H + 1$ Booleans, one per numerator $a = i - H$, all true at the start; for each
  modulus in turn, every entry is replaced by its conjunction with the row of $b$ at $a$. Its
  contract: entry $i$ is the survivor test of the numerator $i - H$. These are the ANDs of the
  first stage of ratpoints in the build with 64-bit bit arrays (\csource{sift.c}{943}{1022}) and
  of its second stage (\csource{sift.c}{1184}{1202}); the default build does the first in chunks
  of registers.
\end{definition}

\begin{definition}[The survivors of one denominator]\label{def:checkDen}
  \uses{def:bruteForce}
  Proposed names: \lname{Ratpoints.Sieve.checkDen}, \lname{Ratpoints.Sieve.checkDen\_spec}.
  The numerators whose entry in a given array of Booleans is true, in increasing order
  (\csource{sift.c}{1213}{1296}), each tested for common factors with $b$
  (\csource{sift.c}{318}{343}, \csource{sift.c}{438}{446}) and checked exactly: $F_c(a, b)$ by
  Horner's rule on the products $c_k b^{d-k}$, computed once per denominator, then $(a, y, b)$
  and, for $y \ne 0$, $(a, -y, b)$ when $F_c(a, b) = y^2$ with $y \ge 0$
  (\csource{find_points.c}{5262}{5304}, \csource{find_points.c}{5312}{5326}); the exact check is
  that of the brute-force reference. Its contract: if entry $i$ is $k(i - H)$ for every $i$, for a
  test $k$ of the numerators, then the result is what the brute-force reference finds over the
  denominator $b$ and the numerators $-H, \dots, H$ that pass $k$.
\end{definition}

\begin{definition}[The denominators]\label{def:sieveRun}
  \uses{def:sieveDen, def:checkDen}
  Proposed names: \lname{Ratpoints.Sieve.sieveRun}, \lname{Ratpoints.Sieve.sieveRun\_spec}.
  The denominators $b = 1, \dots, H$ in increasing order, each sieved and its survivors checked,
  the points appended in turn (\csource{find_points.c}{4059}{4075},
  \csource{find_points.c}{5152}{5159}). Its contract: the result is the concatenation over
  $b = 1, \dots, H$ of what the brute-force reference finds over $b$ and the numerators that
  survive the moduli for $b$, that is, the part of the reference with $z \ge 1$, its numerators
  filtered by the moduli.
\end{definition}

\begin{definition}[The points at infinity]\label{def:pointsAtInfinity}
  \uses{def:bruteForce}
  Proposed names: \lname{Ratpoints.Sieve.pointsAtInfinity},
  \lname{Ratpoints.Sieve.pointsAtInfinity\_eq\_pointsOfDen}.
  $(1, 0, 0)$ for odd $d$; for even $d$, $(1, s, 0)$ and $(1, -s, 0)$ when $c_d = s^2$ with
  $s > 0$, $(1, 0, 0)$ when $c_d = 0$, and nothing when $c_d$ is not a square
  (\csource{find_points.c}{4894}{4899}, \csource{find_points.c}{5068}{5105}). These are the
  points of the candidate $(1, 0)$ of the brute-force reference.
\end{definition}

\begin{definition}[The plain sieve]\label{def:findPoints}
  \uses{def:Params, def:sieveRun, def:pointsAtInfinity, def:maxDegree}
  Proposed names: \lname{Ratpoints.Refusal}, \lname{Ratpoints.findPoints}.
  Refuse more than $\mathrm{maxDegree} + 1$ coefficients, fewer than two, and $H = 0$, in the order
  of the program,
  with its errors 3, 5 and~4 (Section~\ref{sec:refusals}); the library refuses the same inputs
  as bad arguments (\csource{find_points.c}{4557}{4558}, \csource{find_points.c}{4633}{4633},
  \csource{find_points.c}{4658}{4658}). Otherwise make the table of each modulus and return the
  points at infinity followed by the points of the denominators.
\end{definition}

Each component differs from the C code in a stated way, recorded in the \verb|Deviation:|
paragraph of the component whose range holds it:
\begin{itemize}
  \item the Booleans stand for the bit arrays of words, the pieces and the end masks (C2), and the
    walk upwards for the walk by trailing zeros;
  \item every denominator from $1$ to $H$, one after the other, for the loop shapes, the tests on
    the denominators and the blocks (E), and the threads, which deliver the blocks in the same
    order;
  \item every numerator from $-H$ to $H$ for the search region (H) and the 2-adic classes (E, F);
  \item one table per prime and every modulus applied to every numerator before the test for
    common factors, for the composite moduli, the packing, the three stages and their correction
    during the run (F), whose third stage runs after the test for common factors and skips the
    primes that divide $b$;
  \item a plain $\gcd$ for the branch-free binary $\gcd$ of \verb|relprime| (F);
  \item Horner's rule modulo $p$ reduced at every step for the reduction of ratpoints on a fixed
    schedule (F), and the coefficients as integers for its two paths through \verb|long|s and GMP
    integers (D);
  \item the products $c_k b^{d-k}$ computed for every denominator, where ratpoints computes them
    when the first survivor reaches the check;
  \item no reversal (D), and a list of points for the callback \verb|process|, with no stop it may
    request, no \verb|NO_Y| and no \verb|NO_CHECK| (G);
  \item no failure for lack of memory, which the model cannot meet.
\end{itemize}
C1 runs on every input it does not refuse, including the zero polynomial and the inputs that
ratpoints refuses as not squarefree, on which it returns the points of the specification; there
the points at infinity need the case $y = 0$, which ratpoints never meets.

\section{The theorems}\label{sec:sieveTheorems}

\begin{lemma}[Sound filters]\label{lem:flatMap_sound}
  \uses{def:bruteForce}
  Proposed names: \lname{List.flatMap\_filter\_of\_ne},
  \lname{Ratpoints.BruteForce.pointsOfDen\_filter\_of\_ne}.
  If every element $x$ of a list with $g(x)$ nonempty passes the test $k$, then filtering the
  list by $k$ before applying $g$ does not change the concatenation of the $g(x)$. In particular a
  filter of the numerators that every point passes does not change what the brute-force reference
  finds over a denominator.
\end{lemma}

\begin{lemma}[Every point survives]\label{lem:survives}
  \uses{def:admits, def:Params}
  Proposed name: \lname{Ratpoints.Sieve.survives\_of\_mem\_pointsOfDen}.
  For well-formed parameters, a numerator over which the exact check finds a point for the
  denominator $b$ survives the tables of the moduli for $b$.
\end{lemma}

\begin{proof}
  \uses{lem:admits_examinePrime}
  Every modulus is a prime of the table; apply Lemma~\ref{lem:admits_examinePrime} to each.
\end{proof}

\begin{theorem}[The plain sieve returns the reference]\label{thm:findPoints_eq_bruteForce}
  \uses{def:findPoints, def:Params, def:bruteForce}
  Proposed names: \lname{Ratpoints.findPoints\_eq\_ok\_bruteForce},
  \lname{Ratpoints.eq\_bruteForce\_of\_findPoints\_eq\_ok}.
  For well-formed parameters, $2$ to $\mathrm{maxDegree} + 1$ coefficients and $H \ge 1$,
  $\mathrm{findPoints}(c, H, P)$ returns $\mathrm{bruteForce}(c, H)$; so for well-formed
  parameters, every list that $\mathrm{findPoints}$ returns is that of the brute-force reference.
\end{theorem}

\begin{proof}
  \uses{def:sieveRun, def:pointsAtInfinity, lem:survives, lem:flatMap_sound}
  By the contract of the denominators, the sieve returns the part of the reference with
  $z \ge 1$, the numerators of each denominator filtered by the survivor test; by
  Lemma~\ref{lem:survives} every numerator over which the exact check finds a point survives, so
  by Lemma~\ref{lem:flatMap_sound} the filter changes nothing. The points at infinity are those of
  the candidate $(1, 0)$ of the reference.
\end{proof}

\begin{theorem}[The points, each once; the refusals]\label{thm:findPoints}
  \uses{def:findPoints, def:Params, def:IsPoint, def:points}
  Proposed names: \lname{Ratpoints.mem\_iff\_isPoint\_of\_findPoints\_eq\_ok},
  \lname{Ratpoints.nodup\_of\_findPoints\_eq\_ok},
  \lname{Ratpoints.toFinset\_eq\_points\_of\_findPoints\_eq\_ok},
  \lname{Ratpoints.exists\_findPoints\_eq\_error\_iff}.
  For well-formed parameters, if $\mathrm{findPoints}(c, H, P)$ returns the list $L$, then a
  triple lies in $L$ if and only if it is a normalized point of height at most $H$, $L$ has no
  duplicates, and its set of elements is $\mathrm{points}(c, H)$. $\mathrm{findPoints}(c, H, P)$
  refuses the input exactly when $c$ has fewer than two or more than $\mathrm{maxDegree} + 1$
  entries or $H = 0$.
\end{theorem}

\begin{proof}
  \uses{thm:findPoints_eq_bruteForce, thm:bruteForce}
  The first three statements follow from Theorem~\ref{thm:findPoints_eq_bruteForce} and
  Theorem~\ref{thm:bruteForce}; the refusals are the three tests of $\mathrm{findPoints}$.
\end{proof}

\begin{verbatim}
theorem mem_iff_isPoint_of_findPoints_eq_ok (hP : P.WF)
    (h : findPoints c H P = .ok L) : p ∈ L ↔ IsPoint c H p
theorem nodup_of_findPoints_eq_ok (hP : P.WF) (h : findPoints c H P = .ok L) :
    L.Nodup
theorem exists_findPoints_eq_error_iff :
    (∃ e, findPoints c H P = .error e) ↔
      c.length < 2 ∨ maxDegree + 1 < c.length ∨ H = 0
theorem wf_sieveParams (o : SieveOptions) : (sieveParams o).WF
\end{verbatim}

\verb|mem_iff_isPoint_of_findPoints_eq_ok| and \verb|nodup_of_findPoints_eq_ok| are the main
theorems of the project; later stages add the options of \verb|findPoints| as an argument and,
from stage~H on, the hypothesis of Sturm's theorem. The refusal theorem has its C1 form: stage~D
adds the disjunct that $\Phi_c$ is not squarefree. The equation with the brute-force reference
holds until the reversal of stage~D and the divisor order of stage~E change the order of the
output; from then on the theorem is that the output is a permutation of the reference. Comparator
certifies the two main theorems, the refusal theorem, which rules out a \verb|findPoints| that
refuses every input, and \verb|wf_sieveParams|: the challenge states them with \verb|sorry| over
\verb|findPoints|, \verb|Params|, \verb|Params.WF|, \verb|SieveOptions| and \verb|sieveParams|,
imported from the model, never restated, since a definition of the model elaborates without
Mathlib to other instance terms. Each main theorem gets a guard on its axioms. The contracts need
only the program logic of Lean core and stay in a file without Mathlib, so that their
\verb|grind| calls see only the lemmas of core, untouched by Mathlib's updates; that file also
states the result of \verb|sieveRun| as a plain equation, so that the files with Mathlib need no
program logic. The soundness of the rows needs Mathlib, through $\ZZ/p\ZZ$.

\section{The command-line program and the differential tests}\label{sec:sieveCli}

\verb|ratpoints-lean| reads the values of \verb|-p|, \verb|-n|, \verb|-N|, \verb|-P|, \verb|-F|
and \verb|-j|/\verb|-J| into \verb|SieveOptions|, with the checks of \verb|main.c| unchanged, and
runs \verb|findPoints| with \verb|sieveParams| of them; the other options stay as in stage~B. So
the reading of the options (Definition~\ref{def:cliInput}) returns these values from now on, and
the program (Definition~\ref{def:cliRun}) prints the points of the sieve. Since these are those of
the brute-force reference, every differential test of stage~B applies unchanged, and the brute
force needs no test of its own: Theorem~\ref{thm:findPoints_eq_bruteForce} says that the sieve
returns its list. Compiled, the plain sieve handles about $1.5 \cdot 10^7$ candidates per second
with the 30 default moduli: $0.13$~seconds at $H = 1000$, against $0.3$ to $0.9$~seconds for the
reference; so the tests rise to heights of a few thousand, while the runs of \verb|test4.sh| with
heights from $30000$ to $100000$, minutes to tens of minutes each, wait for the bit arrays of
stage~C2. A further suite runs both programs with other parameters (no moduli, all 53 primes,
the cut points in several positions, a third stage) on the same inputs.

\section{Files, names and order of work}

The interface (\verb|primeTable|, \verb|Params|, \verb|SieveOptions|, \verb|sieveParams|,
\verb|Refusal|, \verb|findPoints|) lives in the namespace \verb|Ratpoints|, the components of the
sieve in \verb|Ratpoints.Sieve|, and the general list lemma in \verb|Ratpoints/Mathlib| under its
name in Mathlib's namespace. The table of primes and the constants of the build join
$\mathrm{maxDegree}$ in \verb|Ratpoints/Model/Constants.lean|; the rest of the model goes into
\verb|Ratpoints/Model/Params.lean|,
\verb|Ratpoints/Model/Sieve.lean| and \verb|Ratpoints/Model/FindPoints.lean|; the contracts, without
Mathlib, into \verb|Ratpoints/Proof/Sieve.lean|; the well-formedness and the primes into
\verb|Ratpoints/Proof/Params.lean|, the tables and the soundness of the rows into
\verb|Ratpoints/Proof/Tables.lean|, and the main theorems into
\verb|Ratpoints/Proof/FindPoints.lean|. One commit per item, each passing the project's checks:
\begin{enumerate}
  \item[C1.1] The table of primes, the constants of the choice of the moduli, and the parameters
    (Definitions~\ref{def:primeTable} to~\ref{def:sieveParams}), with their lemmas.
  \item[C1.2] The tables of a prime and the rows (Definitions~\ref{def:squaresMod}
    to~\ref{def:admits}) with \verb|#guard| examples against the tables of ratpoints, and
    Lemmas~\ref{lem:isSquare_div} and~\ref{lem:admits_examinePrime}.
  \item[C1.3] The sieve of one denominator and its contract, with \verb|#guard| examples.
  \item[C1.4] The survivors of one denominator and its contract, with \verb|#guard| examples.
  \item[C1.5] The denominators and their contract.
  \item[C1.6] The points at infinity, \verb|findPoints| with \verb|#guard| examples against the
    C program, Lemmas~\ref{lem:flatMap_sound} and~\ref{lem:survives},
    Theorems~\ref{thm:findPoints_eq_bruteForce} and~\ref{thm:findPoints} with their axiom guards,
    the comparator challenge with its solution and description, and the main results of
    \verb|formalization.yaml|.
  \item[C1.7] \verb|ratpoints-lean| on the sieve, with the options of \verb|SieveOptions|, the
    nodes of the program and the description of the repository brought up to date, and the
    differential tests with their new heights and the suite of parameters.
\end{enumerate}

\section{The evidence}\label{sec:sieve-evidence}

Three checks, made while this chapter was written:
\begin{itemize}
  \item a reading of the C code of a default run at the pin, which fixed the steps of
    Section~\ref{sec:skeleton}, the meaning of a row (the entry of \verb|is_f_square| at
    $a b^{-1}$ for $p \nmid b$, ``$p \nmid a$'' for $p \mid b$, where the curve plays no part),
    and the fact that the points found and their order do not depend on the moduli, their order or
    the cut points;
  \item a prototype in Lean of the model of this chapter and of all its proofs: the model agrees
    with the brute-force reference on examples, among them the zero polynomial, a polynomial of
    degree $0$, no moduli and all 53 primes; the three contracts go through with \verb|vcgen| in
    at most 12 lines each and under half a second each (a rewrite of the entries of the array
    before \verb|grind|, which otherwise times out on the bound proofs that \verb|vcgen| builds for
    \verb|Vector.set|); the soundness of the rows, the equation with the reference and its
    corollaries take the standard axioms only, with no proof longer than 12 lines;
  \item the prototype compiled as a program and timed against the brute-force reference on four
    curves at heights 100, 300 and 1000.
\end{itemize}
